% Physical unit 014; accepted Session 02.
\GordanUnitStart{014}{printed page 9}
\GordanTopFolio{9}
They therefore are equal to:
\[
 {\pi\over 2^{2m+1}}\cdot
 {2m\cdot(2m-1)\cdot(2m-2)\cdots(m+1)\over 1\cdot2\cdot3\cdots m}.
\]
If these values are reduced to our formula, there results:
\[
 {2\over b\pi}\Sigma(l_0)=1+o\sin^2l_0-{3\over4}o^2\sin^4l_0+{5\over4}o^3\sin^6l_0-{175\over64}o^4\sin^8l_0\ldots .
\]

\begin{center}\S\ 3.\end{center}

From this you see that, when
\[
 \Sigma(l_0')>\Sigma(l_0''),
\]
one has
\[
 l_0'>l_0'',
\]
which agrees with the fact that $l_0$ is a singly fitted function of the period of extension, or that a geodetic line is completely determined by its period of extension.

Therefore, by the inverse series and putting
\[
 {2\over b\pi}\cdot\Sigma(l_0)-1=u,
\]
there exists:
\[
 o\sin^2l_0=u+{3\over4}u^2-{1\over8}u^3+{5\over32}u^4.
\]

\begin{center}\S\ 4.\end{center}

To unravel the period of longitude, computing
\[
 {d\Sigma(l_0)\over dl_0}
\]
and using the identity
\[
 {d\Sigma(l_0)\over dl_0}=a\cos l_0\,{d\mathcal A(l_0)\over dl_0},
\]
we obtain
\[
\begin{aligned}
 a\cos l_0\,{d\mathcal A(l_0)\over dl_0}
 &= {b\over2}\cdot\pi\cos l_0\{2o\sin l_0-3o^2\sin^3l_0+{15\over2}o^3\sin^5l_0\\
 &\hspace{33mm}{}-{175\over8}o^4\sin^7l_0\ldots\}\\
 {d\mathcal A(l_0)\over dl_0}
 &= {b\over a}\cdot{\pi\over2}\sin l_0\{2o-3o^2\sin^2l_0+{15\over2}o^3\sin^4l_0\\
 &\hspace{33mm}{}-{175\over8}o^4\sin^6l_0\ldots\}\\
 &= {b\pi\over2a}\sin l_0\{(2-3o+{15\over2}o^2-{175\over8}o^3)o+(3-15o+{525\over8}o^2)o^2\cos^2l_0\\
 &\hspace{35mm}{}+({15\over2}-{525\over8}o)o^3\cos^4l_0+{175\over8}o^4\cos^6l_0\ldots\}.
\end{aligned}
\]
And integrating between the limits $0$ and ${\pi\over2}$,
\[
\begin{aligned}
 {\pi\over2}-\mathcal A(l_0)&={b\over a}{\pi\over2}\{(2-3o+{15\over2}o^2-{175\over8}o^3)o\cos l_0+(1-5o+{175\over8}o^2)o^2\cos^3l_0\\
&\hspace{31mm}{}+({3\over2}-{105\over8}o)o^3\cos^5l_0+{25\over8}o^4\cos^7l_0,\ldots\}.
\end{aligned}
\]
